\section{Orders below thirty-two}
\label{sec:orders-below-thirty-two}

We develop those local certificates and complete the remaining range.
Orders \(4\) through \(9\) follow from low-degree configurations.  For orders
\(10\) through \(31\), a sparse-set argument, supplemented by a short endpoint
census, excludes the shared hub forced by the incidence count.

\subsection{Additional low-order cut certificates}

For a set \(A\subseteq V\), write
\(\partial A=e(A,V\setminus A)\).  We use the following direct
consequences of the variational principle.

\begin{proposition}[Basic cut certificates]\label{lem:cut-certificates}
Let \(G\) have \(n\ge4\) vertices.  Each of the following conditions implies
\(\ac(G)\le2\).
\begin{enumerate}[label=\textup{(\roman*)}]
  \item \(\varnothing\ne A\ne V\) and
  \(n\,\partial A\le2|A|(n-|A|)\).
  \item Disjoint nonempty sets \(A_+,A_-\) have the same cardinality \(k\),
  and, for \(O=V\setminus(A_+\cup A_-)\),
  \(4e(A_+,A_-)+e(A_+,O)+e(A_-,O)\le4k\).
  \item Vertices \(x,y,z\) form a triangle and
  \(n(d(x)+d(y)+d(z)-6)\le6(n-3)\).
  \item Four distinct vertices \(a,b,c,d\) induce exactly the two edges
  \(ab\) and \(cd\), that is, an induced copy of \(2K_2\), and
  \(d(a)+d(b)+d(c)+d(d)\le12\).
\end{enumerate}
\end{proposition}

\begin{proof}
For (i), assign the value \(n-|A|\) on \(A\) and \(-|A|\) on its
complement.  This vector has squared norm \(n|A|(n-|A|)\) and Dirichlet
energy \(n^2\partial A\); hence the hypothesis makes its Rayleigh quotient
\(\le2\), and \eqref{eq:variational} proves (i).  For (ii), use the vector
taking the values \(1,-1,0\) on \(A_+,A_-,O\), respectively.  Its squared
norm is \(2k\), and its Dirichlet energy is the left-hand side of the stated
inequality, so \eqref{eq:variational} likewise proves (ii).  Part (iii)
follows from (i), because the triangle has boundary
\(d(x)+d(y)+d(z)-6\).  For (iv), apply (ii) with \(A_+=\{a,b\}\) and
\(A_-=\{c,d\}\); their total boundary into \(O\) is \(\le12-4=8\).
\end{proof}

The following elementary fact converts the absence of the last two
certificates into a restriction on a low-degree induced subgraph.  It also
follows from the structural results of Chung et
al.~\cite[Corollary to Theorem~1 and Theorem~2]{chung1990}; for completeness,
we give a direct proof for graphs with maximum degree \(\le3\).

\begin{lemma}[Triangle-free induced matching]
\label{lem:triangle-free-induced-matching}
Every triangle-free graph \(J\) on at least eight vertices, with no
isolated vertex and maximum degree \(\le3\), contains an induced
copy of \(2K_2\).
\end{lemma}

\begin{proof}
Suppose otherwise, and choose a vertex \(x\) of maximum degree
\(\Delta\).  Every vertex is at distance \(\le2\) from \(x\).
Indeed, if \(z\) were farther away, choose \(z'\in N(z)\).  The set
\(N(x)\) cannot be contained in \(N(z')\), since
\(z\in N(z')\setminus N(x)\) would give \(d(z')>\Delta\).  An edge
from \(x\) to a vertex in \(N(x)\setminus N(z')\), together with
\(zz'\), would then induce \(2K_2\).

Let \(B\) be the vertices at distance \(2\) from \(x\), and put
\(C_b=N(b)\cap N(x)\) for \(b\in B\).  Each \(C_b\) is nonempty.
Triangle-freeness and the absence of an induced \(2K_2\) imply, for
distinct \(b,b'\in B\),
\[
 C_b\cap C_{b'}=\varnothing\quad\Longleftrightarrow\quad bb'\in E(J),
 \qquad
 bb'\in E(J)\quad\Longrightarrow\quad C_b\cup C_{b'}=N(x).
\]
Every member of \(N(x)\) lies in \(\le\Delta-1\) of the sets
\(C_b\), so \(\sum_{b\in B}|C_b|\le\Delta(\Delta-1)\).  Hence
\(|B|\le2\) if \(\Delta\le2\).  If \(\Delta=3\) and \(|B|\ge4\),
at least two of the \(C_b\) are singletons.  They cannot be distinct,
and if two are equal, their common element already lies in
\(\Delta-1=2\) of the sets \(C_b\).  Every remaining \(C_b\) is therefore
disjoint from that singleton and, by the preceding implications, equals its
complement in \(N(x)\).  A vertex indexed by such a complementary
two-set is adjacent to both singleton occurrences and to two vertices in
\(N(x)\), giving it degree at least \(4\), a contradiction.  Thus
\(|B|\le3\), and in every case
\(|V(J)|\le1+\Delta+|B|\le7\).
\end{proof}

\subsection{Small orders \texorpdfstring{\(4\le n\le9\)}{4 <= n <= 9}}

The preceding cut certificates settle the smallest orders directly.

\begin{proposition}\label{prop:orders-4-9}
If \(4\le n\le9\) and \(G\) has \(n\) vertices and \(2(n-2)\) edges, then
\(\ac(G)\le2\).
\end{proposition}

\begin{proof}
For \(n\le7\), the degree sum satisfies \(4n-8<3n\), so
\(\delta(G)\le2\).  The closed neighborhood of a vertex of degree
\(\le2\) has \(\le3\) vertices, and Lemma~\ref{lem:low-degree}
applies.

It remains to consider \(n=8,9\).  We may again assume
\(\delta(G)\ge3\).  Write \(D_3\) for \(D_3(G)\).  By
\eqref{eq:total-excess},
\(|V\setminus D_3|\le n-8\le1\), so \(|D_3|\ge8\).  Every vertex of
\(G[D_3]\) has degree between \(2\) and \(3\) in \(G[D_3]\).  If
\(G[D_3]\) contains a triangle, its three
vertices all have degree \(3\) in \(G\), and
Proposition~\ref{lem:cut-certificates}(iii) applies.  Otherwise
Lemma~\ref{lem:triangle-free-induced-matching} gives an induced
\(2K_2\) in \(G[D_3]\); its four endpoints have total degree \(12\), so
Proposition~\ref{lem:cut-certificates}(iv) applies.
\end{proof}

\subsection{The shared-hub range \texorpdfstring{\(10\le n\le31\)}{10 <= n <= 31}}

We first show that any obstruction in this range must contain a degree-\(4\)
vertex adjacent to two degree-\(3\) vertices.

\begin{lemma}[Shared-hub existence]\label{lem:shared-hub-existence}
Every degree-\(3\)-separated obstruction of order \(10\le n\le31\)
has a degree-\(4\) vertex with at least two neighbors in \(D_3(G)\).
\end{lemma}

\begin{proof}
Suppose every degree-\(4\) vertex has \(\le1\) such neighbor, and let
\(c\) be the number of vertices of degree at least \(5\).  A degree-\(4\)
vertex contributes \(\le d(x)-3\) to the incidence count, while a
vertex of degree at least \(5\) contributes \(\le(d(x)-3)+3\).
Also \(c\le X\).  Using
\eqref{eq:incidence-total} and
\(\sum_{d(x)\ge4}(d(x)-3)=n-8\) gives
\[
 24+3X
 \le\sum_{d(x)\ge4}(d(x)-3)+3c
 \le n-8+3X,
\]
which forces \(n\ge32\), a contradiction.
\end{proof}

It therefore suffices to exclude this shared-hub configuration.  The
general sparse-set argument reduces its only difficult branch to orders
\(10\) through \(15\); the following elementary incidence fact is needed
at that endpoint.

\begin{lemma}[Six-column incidence lemma]\label{lem:six-column}
Let \(A_0,A_1,A_2\) be disjoint two-element sets and let \(B\) have six
elements.  Suppose \(J\) is a bipartite graph with parts
\(A_0\cup A_1\cup A_2\) and \(B\), and every vertex of \(B\) and both
vertices of \(A_0\) have degree \(3\) in \(J\).  Then two vertices of
\(A_0\cup A_1\cup A_2\) have at least two common neighbors in \(B\),
where either they belong to the same \(A_i\), or one belongs to \(A_0\)
and the other to \(A_1\cup A_2\).
\end{lemma}

\begin{proof}
Suppose not.  For \(b\in B\), put
\(\alpha_b=|N_J(b)\cap A_0|\), \(\beta_b=|N_J(b)\cap A_1|\), and
\(\gamma_b=|N_J(b)\cap A_2|\), and regard
\((\alpha_b,\beta_b,\gamma_b)\) as the column indexed by \(b\).  For a
condition \(P\), let \(\mathbf1_P\) denote its indicator.  Then
\(\alpha_b,\beta_b,\gamma_b\le2\), and
\(\alpha_b+\beta_b+\gamma_b=3\).  Moreover,
\[
 \sum_{b\in B}\alpha_b=6,\qquad
 \sum_{b\in B}\!\bigl(\mathbf1_{\alpha_b=2}
   +\mathbf1_{\beta_b=2}+\mathbf1_{\gamma_b=2}\bigr)\le3,
\]
while
\(\sum_b\alpha_b\beta_b\le4\) and
\(\sum_b\alpha_b\gamma_b\le4\).

Call the column indexed by \(b\) \emph{balanced} if
\((\alpha_b,\beta_b,\gamma_b)=(1,1,1)\), and let \(r_{\mathrm{bal}}\) be the
number of balanced columns.  Every other column has an entry equal to \(2\),
so \(r_{\mathrm{bal}}\ge3\).  Pointwise,
\[
 \mathbf1_{\{b\ {\rm balanced}\}}+2\mathbf1_{\{\gamma_b=0\}}
 \le\alpha_b\beta_b.
\]
After summing, \(r_{\mathrm{bal}}+2|\{b:\gamma_b=0\}|\le4\), and hence
\(\gamma_b\ge1\) for every \(b\).  The same argument with
\(\beta\) and \(\gamma\) interchanged gives \(\beta_b\ge1\).
Consequently \(\alpha_b\le1\), with equality only in a balanced
column.  Thus \(6=\sum_b\alpha_b\le r_{\mathrm{bal}}\), so all six columns are
balanced.  But then \(\sum_b\alpha_b\beta_b=6\), contradicting its
upper bound \(4\).
\end{proof}

We now isolate the finite endpoint of the shared-hub argument.  The
statement includes only the hard branch; the two easy branches will be
handled uniformly afterward.

\begin{proposition}[Endpoint shared-hub census]
\label{prop:endpoint-shared-hub}
Let \(10\le n\le15\), let \(G\) have \(n\) vertices and \(2(n-2)\)
edges, and suppose that \(\delta(G)\ge3\) and \(D_3(G)\) is independent.
Let \(z\) have degree \(4\) and distinct neighbors \(u,v\in D_3(G)\).
Put \(A=\{z,u,v\}\), \(F=N(A)\setminus A\), and
\(B=V\setminus(A\cup F)\).
If \(e(B,F)>2|B|\) and \(I(B)\le2\varepsilon(B)\), then \(\ac(G)\le2\).
\end{proposition}

\begin{proof}
Write \(f=|F|\), \(b=|B|\), \(q=e(B,F)\),
\(I_B=I(B)=2e(B)\), \(I_F=I(F)=2e(F)\), and
\(\varepsilon_F=\varepsilon(F)\), \(\varepsilon_B=\varepsilon(B)\).  Set
\(F_0=N(z)\setminus A\), \(F_1=N(u)\setminus A\), and
\(F_2=N(v)\setminus A\).  These sets all have two elements, and
\(F=F_0\cup F_1\cup F_2\); hence
\(f\le6\).  Every vertex of \(K=F_1\cup F_2\) has degree at least
\(4\), because \(D_3(G)\) is independent.  In particular,
\begin{equation}\label{eq:endpoint-k-excess}
  |K|\le \varepsilon_F.
\end{equation}

The excess on \(A\) is \(1\), and exactly six edges join \(A\) to \(F\).
The corresponding degree and incidence identities are
\begin{align}
 \varepsilon_F+\varepsilon_B&=n-9,                 \label{eq:endpoint-excess}\\
 q+I_B&=3b+\varepsilon_B,                           \label{eq:endpoint-bulk}\\
 I_F+q&=3f+\varepsilon_F-6.                         \label{eq:endpoint-moat}
\end{align}
Also \(q\le \varepsilon_F+2f\).  Combining this bound with
\eqref{eq:endpoint-bulk} and \(I_B\le2\varepsilon_B\) gives
\(3b\le n-9+2f\).  Since \(b=n-3-f\),
\begin{equation}\label{eq:endpoint-f}
  2n\le5f.
\end{equation}

For \(n=10\), \eqref{eq:endpoint-excess} gives
\(\varepsilon_F+\varepsilon_B=1\), contrary to \(|K|\ge2\) and
\eqref{eq:endpoint-k-excess}.  For \(n=11\),
\eqref{eq:endpoint-f} gives \(f\ge5\), so
\(|K|\ge f-|F_0|\ge3>\varepsilon_F+\varepsilon_B=2\), again a contradiction.

For the four remaining orders, equations
\eqref{eq:endpoint-excess}--\eqref{eq:endpoint-f}, the evenness of
\(I_B\), and \eqref{eq:endpoint-k-excess} give precisely the following
profiles.  A tuple records
\((f,b;\varepsilon_F,\varepsilon_B;I_B,q,I_F)\):
\[
\begin{array}{c|c}
n&\text{possible profiles}\\ \hline
12&(5,4;3,0;0,12,0)\\
13&(6,4;4,0;0,12,4)\\
14&(6,5;5,0;0,15,2),\
    (6,5;4,1;0,16,0),\
    (6,5;4,1;2,14,2)\\
15&(6,6;6-j,j;2j,18-j,0),\quad j=0,1,2.
\end{array}
\]
These profiles follow directly from the preceding constraints.  At \(n=12\),
\eqref{eq:endpoint-f} gives \(f\ge5\); if \(f=6\), then
\(|K|\ge4>\varepsilon_F+\varepsilon_B=3\).  Thus \(f=5\), after which
\eqref{eq:endpoint-k-excess} forces \(\varepsilon_F=|K|=3\) and
\(\varepsilon_B=0\).
At \(n=13,14,15\), equation \eqref{eq:endpoint-f} forces \(f=6\).
The three pairs \(F_i\) are then disjoint, so \(|K|=4\) and
\(\varepsilon_B\le n-13\).  Substitution in
\eqref{eq:endpoint-bulk} and \eqref{eq:endpoint-moat}, together with
\(0\le I_B\le2\varepsilon_B\), gives the displayed rows.

It remains to eliminate the profiles.

\smallskip
\noindent\emph{Orders \(12\) and \(13\).}
In both rows, every vertex of \(B\) has degree \(3\).  Moreover,
\(\varepsilon_F=|K|\); since every vertex of \(K\) contributes at least one
unit to \(\varepsilon_F\), no excess remains on \(F\setminus K\).  For
\(n=12\), the equalities
\(|F_0|=2\), \(|K|=3\), and \(|F|=5\) also give \(F_0\cap K=\varnothing\);
for \(n=13\), this follows from the pairwise disjointness noted above.
Hence both vertices of \(F_0\) have degree
\(3\).  They cannot be adjacent to \(B\), so, besides their common
neighbor \(z\), each has two neighbors in \(F\).  This is impossible
when \(n=12\), where \(I_F=0\).  When \(n=13\), two distinct vertices
of internal degree \(2\) require at least three edges in \(G[F]\), and
hence \(I_F\ge6\), contrary to \(I_F=4\).

\smallskip
\noindent\emph{Order \(14\).}
In the first profile, at least one vertex of \(F_0\) has degree \(3\);
it has two neighbors in \(F\), so \(I_F\ge4\), contrary to \(I_F=2\).
In the second profile, both vertices of \(F_0\) have degree \(3\) and
\(I_F=0\).  Each would need two distinct neighbors of degree at least
\(4\) in \(B\), although \(\varepsilon_B=1\) implies
\(|\{x\in B:d(x)\ge4\}|\le1\).
In the last profile, both vertices of \(F_0\) again have degree \(3\),
and \(|\{x\in B:d(x)\ge4\}|\le1\).  Each vertex
of \(F_0\) must therefore meet \(F\).  Since \(I_F=2\), the unique edge
of \(G[F]\) joins the two vertices of \(F_0\).  Together with \(z\)
they form a triangle of degree sum \(4+3+3=10\), and
  Proposition~\ref{lem:cut-certificates}(iii) applies.

\smallskip
\noindent\emph{Order \(15\).}
Here \(I_F=0\), so every vertex of \(F\) has one neighbor in \(A\) and
all its remaining neighbors in \(B\).  If \(\varepsilon_B=1\), then at least one
vertex of \(F_0\) has degree \(3\) and needs two distinct high-degree
neighbors in \(B\), whereas \(\varepsilon_B=1\) permits only one.

If \(\varepsilon_B=2\), both vertices of \(F_0\) have degree \(3\).  The set
\(B_+=\{x\in B:d(x)\ge4\}\) consequently consists of two degree-\(4\)
vertices, and both vertices of \(F_0\) are adjacent to both members of
\(B_+\).  For \(Q=\{z\}\cup F_0\cup B_+\), the five vertices, listed with
\(z\) first, have respectively \(\le2,0,0,2,2\) neighbors outside \(Q\).
Thus \(\partial Q\le6\).
Since \(15\cdot6\le2\cdot5\cdot10\),
Proposition~\ref{lem:cut-certificates}(i) applies.

Finally, suppose \(\varepsilon_B=0\).  Every vertex of \(B\) has degree \(3\).
The absence of edges in \(F\), together with independence of \(D_3(G)\),
forces every vertex of \(F\) to have degree \(4\).  Thus the bipartite
graph between \(F\) and \(B\) is \(3\)-regular on both sides.
Apply Lemma~\ref{lem:six-column} to the three pairs
\(F_0,F_1,F_2\), and denote the two common neighbors in \(B\) by
\(b_1,b_2\).  If the pair \(x,y\) supplied by the lemma lies in one \(F_i\), let
\(p\in A\) be the vertex adjacent to both.  In the order
\((p,x,y,b_1,b_2)\), these vertices have respectively
\(\le2,1,1,1,1\) neighbors outside their set, so
\(\partial\{p,x,y,b_1,b_2\}\le6\).  If instead \(x\in F_0\) and
\(y\in F_1\cup F_2\), let \(p,q\in A\) be adjacent to \(x,y\), respectively.
The vertices \(p,q\) are adjacent, and in the order
\((p,q,x,y,b_1,b_2)\), the six vertices have respectively
\(\le2,1,1,1,1,1\) neighbors outside their set, giving
\(\partial\{p,q,x,y,b_1,b_2\}\le7\).  In the two cases,
\(15\cdot6\le2\cdot5\cdot10\) and
\(15\cdot7\le2\cdot6\cdot9\), respectively, so
Proposition~\ref{lem:cut-certificates}(i) completes the proof.
\end{proof}

\begin{lemma}[Shared degree-\(3\) hub]\label{lem:sharp-shared-hub}
Let \(n\ge10\), let \(G\) have \(n\) vertices and \(2(n-2)\) edges,
and suppose that \(\delta(G)\ge3\) and \(D_3(G)\) is independent.
If a degree-\(4\) vertex has two distinct neighbors in \(D_3(G)\), then
\(\ac(G)\le2\).
\end{lemma}

\begin{proof}
Let the hub be \(z\), let \(u,v\) be the two degree-\(3\) neighbors,
and define \(A,F,B,f,b,q,\varepsilon_F,\varepsilon_B\), and \(I_B\) as in
Proposition~\ref{prop:endpoint-shared-hub}.  The set \(A\) is
two-sparse, \(f\le6\), \(B\ne\varnothing\), and there is no edge from
\(A\) to \(B\).  The same identities give
\[
 q\le \varepsilon_F+2f,\qquad q+I_B=3b+\varepsilon_B,\qquad
 \varepsilon_F+\varepsilon_B=n-9.
\]
If \(q\le2b\), then \(e(A,F)=6=2|A|\), and the two-cluster
certificate applies.  If \(I_B>2\varepsilon_B\),
Proposition~\ref{prop:sparse-set-principle}(ii) produces a two-sparse
nonempty subset of \(B\); together with \(A\), part~(i) gives the
contradiction.

We may therefore assume \(q>2b\) and \(I_B\le2\varepsilon_B\).
As in \eqref{eq:endpoint-f}, these inequalities imply
\(2n\le5f\le30\).  Thus \(n\le15\), and
Proposition~\ref{prop:endpoint-shared-hub} completes the proof.
\end{proof}

\begin{theorem}[Orders below thirty-two]\label{thm:range-4-31}
If \(4\le n\le31\) and \(G\) is a graph on \(n\) vertices with
\(2(n-2)\) edges, then \(\ac(G)\le2\).
\end{theorem}

\begin{proof}
Proposition~\ref{prop:orders-4-9} handles \(n\le9\).  Suppose
\(10\le n\le31\).  If \(G\) has a vertex of degree \(\le2\),
Lemma~\ref{lem:low-degree} applies.  We may therefore assume
\(\delta(G)\ge3\).  If two degree-\(3\) vertices are adjacent, use
Lemma~\ref{lem:adjacent-three}.  Otherwise, a counterexample would be a
degree-\(3\)-separated obstruction.  Lemma~\ref{lem:shared-hub-existence}
produces a degree-\(4\) vertex with two degree-\(3\) neighbors, contrary to
Lemma~\ref{lem:sharp-shared-hub}.
\end{proof}

The ranges established in Sections~\ref{sec:obstruction-closure}
and~\ref{sec:orders-below-thirty-two} are disjoint and exhaustive.

\begin{proof}[Proof of Theorem~\ref{thm:main}]
Apply Theorem~\ref{thm:range-4-31} for \(4\le n\le31\) and
Theorem~\ref{thm:range-ge32} for \(n\ge32\).
\end{proof}
